\documentclass{article}
\usepackage{listings}

\title{在VLC播放器中使用anime4k等glsl滤镜}
\author{youcai}
\date{2025年04月30日 星期三}

\begin{document}

众所周知，vlc的可定制性和播放质量一直不如mpv之类的播放器硬核。尤其是mpv的glsl shader功能是无法在vlc 3.0.x中使用的。由于种种原因，多数动漫的原始分辨率在1920x1080。由于动漫本身的画面特性，vlc默认的二次插值在4k屏幕上看起来十分糟糕，甚至不如1080p点对点显示。

一个显而易见的解决方法是直接用mpv看动漫就完了，然而作为没有gui的硬核播放器，用起来还是多少有点难受的。虽然身为linux老炮，我并不排斥用键盘命令行解决一切，不过看动漫本身就是为了放松一下，还是有个gui方便一点。

但是！随着4.x版本vlc的到来，这一切即将成为历史。新版vlc新增了libplacebo作为可选的画面输出模块，libplacebo不但有一堆的内置gpu加速重采样，甚至支持mpv同款的glsl shader！

\section{安装支持libplacebo的vlc}

目前只要是3.0.x版本的vlc都是不带libplacebo模块的。由于是新功能可能经常有bug，所以建议直接live at head用git主线最新版完事。

\subsection{使用Gentoo}

首先，直接将vlc新版本的封印解除：

\begin{lstlisting}[language=bash]
echo "media-video/vlc **" > /etc/portage/package.accept_keywords/20vlc
\end{lstlisting}

然后，给vlc和ffmpeg打开libplacebo的USE flag：

\begin{lstlisting}[language=bash]
echo "media-video/ffmpeg libplacebo" > /etc/portage/package.use/20vlc
echo "media-video/vlc libplacebo" > /etc/portage/package.use/20vlc
\end{lstlisting}

当然，你也可以直接把libplacebo加到global USE。也需要按需求开其他的USE flag，例如加载ass字幕要打开libass。新版的qt前端也需要qml支持。

装就完了：

\begin{lstlisting}[language=bash]
emerge -avuDN media-video/vlc-9999
\end{lstlisting}

\subsection{使用其他发行版}

这个就～麻烦了。众所周知vlc的libplacebo支持还需要ffmpeg带着libplacebo支持编译。如果你的发行版ffmpeg没开libplacebo，那么你就得先手动编译一个ffmpeg，再编译一个使用这一份ffmpeg的vlc出来。对vlc来说需要带上\verb|--enable-libplacebo|就可以了。这里就不详细的描述过程了～

如果是ubuntu用户，可以使用snap里面的vlc 4.0包。也可以考虑使用gentoo prefix之类。

\section{开启libplacebo输出}

运行VLC之后，按下Ctrl+P快捷键打开设置，会发现左下角多了一个模式切换，点击最右边切换到“专家”模式。

搜索找到\verb|core.vout|，这个设置本来的值应该是\verb|any|，更改成\verb|libplacebo|

\includegraphics{/images/vlc-placebo-1.webp}

\section{开启glsl滤镜}

这里以anime4k为例，首先克隆仓库：

\begin{lstlisting}[language=bash]
git clone https://github.com/bloc97/Anime4K.git
\end{lstlisting}

然后，还是在专家模式设置中，查找这个设置：\verb|placebo.pl-user-shader|

\includegraphics{/images/vlc-placebo-2.webp}

双击编辑，找到刚才克隆的仓库中的glsl shader即可。如果gpu够强，可以直接上超大杯：

\begin{verbatim}
Anime4K/glsl/Upscale+Denoise/Anime4K_Upscale_Denoise_CNN_x2_VL.glsl
\end{verbatim}

退出重启vlc，就可以使用Anime4k了！

\section{运行结果}

使用\verb|vlc -vvv|启动并打开任意视频，找到大致如下的日志，表明成功使用了libplacebo：

\begin{verbatim}
placebo generic debug: Initialized libplacebo v7.349.0 (API v349)
\end{verbatim}

下面不远处应该有这样一行日志，表明user shader加载：

\begin{verbatim}
placebo generic debug: Loaded user shader:
\end{verbatim}

下面会跟着一堆的输出，就是shader的源码。

我的4060ti在使用超大杯扩大1080p动漫到4k的情况下，gpu利用率为\verb|25%|左右，用去611M显存。

\includegraphics{/images/vlc-anime4k.webp}

\end{document}

